\section{从原始披露到投资结论的复算链}

本页选取最容易在研究中发生口径漂移的五项数字，展示原始事实、House 假设、公式与读者可得到的结论。它不是新增模型，而是对“数据--假设--目标价”链条的交叉检查。每次正式业绩披露后，应以同一结构更新，而不能只替换摘要中的利润数字。

\begin{exhibitbox}[表 A-11：五项关键数字的复算实例]
\begin{tabularx}{\exhibitboxwidth}{L{2.5cm}L{3.1cm}X X}
\toprule
项目 & 事实或假设 & 公式 / 勾稽 & 投资含义 \\
\midrule
市值 & 17.75 元；11.5256252 亿股 & $17.75\times11.5256252=204.58$ 亿元 & 所有 PE、PB、EV 的价格日一致 \\
2026E 基准 EPS & 归母 20.10 亿元 & $20.10\div11.5256252=1.7439$ 元 & 现价对应 10.18x PE \\
基准 PE 法价值 & 2026E 1.7439 元；2027E 1.2374 元 & $70\%\times1.7439\times10.5+30\%\times1.2374\times12=17.27$ 元 & 避免只用 2026 高点 \\
基准 PB/ROE 价值 & Q1 BPS 9.572 元；EPS 1.7439 元 & $(9.572+1.7439-0.08)\times1.55=17.42$ 元 & 资产负债表交叉检查 \\
最终目标 & 内在价值 16.95 元；市场锚 17.75 元 & $16.95\times85\%+17.75\times15\%=17.07$ 元 & 展示为 17.0 元；外部目标权重为零 \\
\bottomrule
\end{tabularx}
\end{exhibitbox}

\section{复算通过不等于经营验证通过}

上述算术在当前价格日、股本、概率和四舍五入规则下通过；它不能代替正式 H1 对销量、实现收入代理、单位成本、应收、存货和经营现金流的验证。故模型的正确读法是：17.0 元是一个透明、可被新披露推翻的 12 个月价值判断，而非对短期股价的精确预测。

\sourcenote{验证后的市场价格与股本、财务与预告数据，以及估值复算工作底稿。}
